The "overfitted brain" hypothesis provides a powerful framework for understanding \textit{why} different animals dream differently. If dreams function to "de-overfit" the brain and generalize learning, then the \textit{content} of those dreams should be specific to whatever skills and senses are most critical—and most over-trained—for that species.


\begin{itemize}
    \item \textbf{Rodents and Spatial Replay:} As noted in Section 3.1, rodents replay hippocampal place-cell activity from their waking explorations, effectively re-running mazes in their sleep \citep{wilson1994reactivation}. This aligns perfectly with the hypothesis, as spatial navigation is a critical, high-stakes skill for a rodent, and "replaying" these routes—even in bursts or in reverse—would serve to consolidate and generalize that spatial memory.
    \item \textbf{Bats and Acoustic-Spatial Replay:} This principle extends to other mammals with unique sensory specializations. Bats, for example, possess standard mammalian REM/NREM sleep cycles. Studies have shown that their hippocampus also "replays" flight trajectories during sleep \citep{omer2025batreplay}. Given that their most critical and computationally demanding skill is echolocation, it is inferred that their dreams are likely "echo-spatial remixes" of recent flights. The "bizarre" or "odd reverberations" in such dreams would serve the "overfitting" function perfectly, acting as adversarial noise to help the bat generalize its echolocation skills to new, unseen environments \citep{suga2003batmod}.
    \item \textbf{Insects and Simpler Organisms:} The question of dreaming is meaningful even in simpler animals. Sleep has been clearly identified in insects like \textit{Drosophila} (fruit flies) \citep{tian2023drosophila} and even in nematode worms \citep{lesku2016evolutionary}. These organisms also require sleep for proper brain function and memory consolidation, suggesting the "offline" maintenance function is deeply conserved \citep{tian2023drosophila}. While we cannot know their subjective experience, the core biological need for an offline consolidation phase exists.
    \item \textbf{A Marine Mammal Exception (Dolphins):} The universality of \textit{REM-like} dreaming, however, has a critical exception: cetaceans (dolphins and whales) \citep{lyamin2021cetacean}. These marine mammals have evolved \textbf{unihemispheric slow-wave sleep (USWS)}, a state where one half of the brain sleeps while the other remains "awake." This adaptation is likely necessary to manage breathing in the water and maintain vigilance against predators \citep{ridgway2002unihemispheric, rattenborg2000usws}. Consequently, cetaceans exhibit "little to no REM sleep," as the total-body paralysis and sensory disconnection of REM sleep would be incompatible with their aquatic needs \citep{lyamin2021cetacean}. This suggests that while sleep itself is universal, the \textit{architecture} of sleep (and the associated phenomenology of immersive dreaming) can be radically adapted or suppressed by extreme evolutionary pressures.
\end{itemize}

\begin{tcolorbox}[title={Dream Modalities Across the Tree of Life}, colback=gray!3, colframe=black]
\small
\textit{If dreams evolved to fight overfitting, their “feel” should mirror each species’ dominant learning channel — and their vividness should scale with how much REM-like, bilaterally integrated sleep they can afford.}

\begin{tabular}{@{} p{0.20\linewidth} p{0.23\linewidth} p{0.35\linewidth} p{0.15\linewidth} @{}}
\toprule
\textbf{Taxon} & \textbf{Dominant replay/skill} & \textbf{Likely dream phenomenology (function)} & \textbf{Sleep architecture cue} \\
\midrule
Songbirds (zebra finch) & Song circuits & Auditory “practice”; noisy, re-sequenced motifs — generalize song \citep{dave2000,deregnaucourt2005}. & Robust REM/NREM \\
Rodents (rats/mice) & Spatial navigation (hippocampus) & Map-like route fragments; time-compressed/reverse paths — consolidate routes \citep{wilson1994reactivation,ji2007}. & NREM replay + REM \\
Dogs & Olfaction + social/motor & Smell-rich scenes with motor episodes — update scent maps \& routines \citep{kis2017,rasch2007}. & Mammalian REM present \\
Cats & Predatory motor programs & Action scripts when REM atonia fails (“acting-out”) \citep{mahowald2005}. & Strong REM \\
Bats & Spatial–acoustic flight & Echo-spatial remixes of recent flight paths; odd reverberations — generalize echolocation \citep{omer2025batreplay,suga2003batmod}. & Standard mammalian REM/NREM \\
Dolphins (cetaceans) & Vigilance/breathing & Lateralized, fragmentary quiescence; little/no cinematic content \citep{lyamin2021cetacean,ridgway2002unihemispheric,rattenborg2000usws}. & Unihemispheric SWS; minimal/absent REM \\
Pinnipeds (elephant seals) & Navigation while diving & Short, opportunistic bouts; muted but present REM-like phases \citep{kendallbar2023}. & $\sim$2 h total sleep/day \\
Octopuses & Visual/body patterning & Vivid, rapid skin-display sequences; “active sleep” akin to REM \citep{medeiros2023octopus,ramirez2023octopus,medeiros2021iscience}. & Two-stage quiet/active sleep \\
Jumping spiders & Eye/retinal activity & REM-like bouts with saccades; content unknown \citep{roessler2022pnas}. & REM-like physiology \\
Insects (\textit{Drosophila}) & Learned associations & Circuit reactivation; phenomenology unknown (consolidation) \citep{li2009science,vanalphen2023}. & Sleep homeostasis present \\
\bottomrule
\end{tabular}
\end{tcolorbox}

\textit{[Placeholder: Cautious criteria for inferring “dream-like” across species.]}
